\documentclass{article}
\usepackage{import}
\import{../../template/}{article-template}
\graphicspath{{../../images/}}
\addbibresource{../../template/library.bib}

\title{Wolfram}
\author{PENG}
\date{\today}

\begin{document}
\maketitle
\tableofcontents
\section{Introduction}
\url{https://www.wolfram.com/language/elementary-introduction/}

\begin{definition}[Wolfram Langauge]
    Wolfram Language is a computation language.
\end{definition}

\subsection{Functions}
\begin{definition}[function]~
    \begin{verbatim}
(* Func[params] *)
(* [] for function, () for grouping *)

Plus[x, y, ...]
Times[x, y, ...]
Subtract[x, y] === x + (-1 * y)
Divide[x, y] === x * y^(-1)

Power[x, y, z] === Power[x, Power[y, z]]
x^y^z === x^(y^z)
(x^y)^z === x^(y*z)

Max[1, 2, ...]
Min[1, 2, ...]

RandomInteger[10]
\end{verbatim}
\end{definition}

\subsection{Lists}
\begin{definition}[list]
    Collect things together.
    \begin{verbatim}

\end{verbatim}
\end{definition}

\appendix
\section{Tools}
\begin{itemize}
    \item Palettes $\to$ Code Formatting
    \item \verb|Information[expr, "Properties"]|
    \item directory
          \begin{verbatim}
If[SameQ[$InputFileName, ""],
    If[$Notebooks,
        NotebookDirectory[]
        ,
        Directory[]
    ]
    ,
    DirectoryName[$InputFileName]
]
\end{verbatim}
    \item plot
\end{itemize}

\printbibliography
\end{document}